\documentclass[11pt,a4paper]{article}
\usepackage[T1]{fontenc}
\usepackage[utf8]{inputenc}
\usepackage{lmodern,amsmath,amssymb,amsthm,mathtools,microtype}
\usepackage[margin=29mm]{geometry}
\usepackage[colorlinks=true,linkcolor=blue,citecolor=blue,urlcolor=blue]{hyperref}
\newtheorem{theorem}{Theorem}[section]
\newtheorem{proposition}[theorem]{Proposition}
\newtheorem{lemma}[theorem]{Lemma}
\theoremstyle{remark}
\newtheorem{remark}[theorem]{Remark}
\newcommand{\R}{\mathbb R}
\newcommand{\T}{\mathbb T}
\newcommand{\K}{\mathcal K}
\newcommand{\F}{\mathcal F}
\newcommand{\cg}{c_\gamma}
\newcommand{\dd}{\,\mathrm d}
\newcommand{\norm}[1]{\left\lVert#1\right\rVert}
\title{An extreme solitary wave for the bidirectional Whitham system}
\author{Kaiwen Jin\\\small\texttt{yui.ui.siki@gmail.com}}
\date{15 October 2026}
\hypersetup{pdftitle={An extreme solitary wave for the bidirectional Whitham system},pdfauthor={Kaiwen Jin}}
\begin{document}
\maketitle

\begin{abstract}
We construct an extreme solitary wave for the bidirectional Whitham
system. Its velocity profile is positive, even, continuous, smooth away
from its crest and strictly decreasing to zero on the positive half-line.
The crest reaches the first critical value of the cubic traveling-wave
equation. The construction starts with positive highest periodic waves
of arbitrarily large period. A common modulus of continuity and a fixed
drop from the critical height keep their limit nonconstant. A balance
identity using the kernel's first absolute moment then excludes every
positive limiting pedestal, without assuming an integrable profile.
\end{abstract}

\section{Introduction}\label{sec:intro}

The bidirectional Whitham system is the nonlocal dispersive model
\begin{equation}\label{eq:system}
 \partial_t\eta+\partial_x(K*v+\eta v)=0,\qquad
 \partial_tv+\partial_x\left(\eta+\frac12v^2\right)=0,
\end{equation}
where $\eta$ is the surface elevation, $v$ is the velocity and the
convolution kernel on $\R$ has symbol
\begin{equation}\label{eq:symbol-intro}
 \widehat K(\xi)=\frac{\tanh\xi}{\xi},\qquad \widehat K(0)=1.
\end{equation}
It combines the finite-depth gravity-wave dispersion relation with a
quadratic shallow-water nonlinearity; see \cite{EJC,NW}. We study the
largest crest allowed by its scalar traveling-wave equation.

For profiles moving at speed $c$ and tending to zero at infinity,
integration of \eqref{eq:system} gives $\eta=cv-v^2/2$ and
\begin{equation}\label{eq:profile-intro}
 K*v=c^2v-\frac{3c}{2}v^2+\frac12v^3.
\end{equation}
The derivative of the cubic first vanishes at
$v=c(1-1/\sqrt3)$. At this height the smooth inversion used for smaller
waves degenerates. The existence question is therefore whether a
localized wave can attain this crest while retaining positive, monotone
tails.

Small solitary waves for this system are supplied by the variational
construction of Nilsson and Wang \cite[Theorem 3.1]{NW}. Highest periodic
waves were constructed by Ehrnstr\"om, Johnson and Claassen \cite{EJC};
their Appendix B discusses positive highest periodic waves at the fixed
period of that paper. Ehrnstr\"om, M{\ae}hlen and Varholm
\cite[Section 4.2]{EMV} determine precise crest behavior for extreme
solutions and explicitly leave existence of a highest solitary wave
open in Remark 4.8. A periodic highest wave alone does not answer that
question: as its period grows, it might approach a constant, or retain
a positive value at infinity.

The following result supplies a solitary profile at the critical crest.
\begin{theorem}\label{thm:main}
There exist
\[
 1<c<\sqrt{3(\sqrt3-1)}
\]
and an even continuous function $v:\R\to(0,\infty)$, smooth on
$\R\setminus\{0\}$, such that $v$ is strictly decreasing on
$(0,\infty)$, $\lim_{|x|\to\infty}v(x)=0$, and
\[
 K*v=c^2v-\frac{3c}{2}v^2+\frac12v^3
 \quad\hbox{on }\R,\qquad
 v(0)=c\left(1-\frac1{\sqrt3}\right).
\]
The surface profile $\eta=cv-v^2/2$ satisfies $\eta(0)=c^2/3$.
The pair $(\eta(x-ct),v(x-ct))$ solves \eqref{eq:system} in
distributions, and classically away from the moving crest.
\end{theorem}

The theorem gives an affirmative answer to the existence question in
\cite[Remark 4.8]{EMV}, for the exact symbol and conservative system
\eqref{eq:system}. It is an existence result for one limiting speed;
uniqueness, stability and a family indexed by all speeds are separate
questions.

We obtain the profile as a limit of highest positive periodic waves.
Periodic-to-solitary constructions are available for the unidirectional
Whitham equation; Ehrnstr\"om, Nik and Walker \cite{ENW} give a direct
construction of a full family. That equation has a different multiplier
and a quadratic profile nonlinearity, and its treatment of nonzero tails
uses a Galilean transformation. Here the cubic retains the speed in
several coefficients. We instead integrate the convolution defect over
expanding symmetric intervals. The first absolute moment of $K$ makes
this defect tend to zero for any bounded profile with equal tails. Strict
concavity of the cubic contradicts that balance if the limiting tail is
positive.

Two estimates allow the period to grow. A kernel-translation estimate
gives one modulus of continuity for all heights up to the critical value,
and a differentiated equation for smooth periodic waves gives a drop from
that height on a fixed interval. Taking highest limits preserves the drop
without differentiating at a crest. We also construct the positive
periodic branch for every sufficiently large period, using the corrected
analytic continuation theorem of \cite{CSV}; this step is needed beyond
the fixed-period positive-wave discussion in \cite{EJC}.

Section~\ref{sec:estimates} gives the kernel facts and the uniform
estimates. Section~\ref{sec:periodic} constructs the positive highest
periodic waves, and Section~\ref{sec:solitary} passes to the real line and
proves the theorem. Throughout,
$\widehat f(\xi)=\int_\R e^{-ix\xi}f(x)\dd x$, with inverse factor
$1/(2\pi)$, and convolution uses Lebesgue measure.

\section{The kernel and estimates uniform in the period}\label{sec:estimates}

Write
\begin{equation}\label{eq:constants}
 a=1-\frac1{\sqrt3},\qquad \gamma(c)=ac,\qquad
 F_c(t)=c^2t-\frac{3c}{2}t^2+\frac12t^3,
 \qquad \cg=\sqrt{3(\sqrt3-1)}.
\end{equation}
The first positive critical point of the cubic is $\gamma(c)$. In fact,
\begin{equation}\label{eq:cubic-derivatives}
 F_c'(t)=\frac32(\gamma(c)-t)
       \left(\gamma(c)-t+\frac{2c}{\sqrt3}\right),\qquad
 F_c''(t)=3(t-c).
\end{equation}
Thus $F_c$ is strictly increasing and strictly concave on
$[0,\gamma(c)]$ when $c>0$. A profile with maximum $\gamma(c)$ will be
called highest. Smooth profiles lying strictly below this value will be
called subcritical; this refers to their height, rather than their speed.

The normalized kernel is
\begin{equation}\label{eq:kernel}
 K(x)=\frac1\pi\log\coth\frac{\pi|x|}{4},\qquad x\ne0.
\end{equation}
This is the formula in \cite[Section 4.2]{EMV} for the symbol used here.
It fixes the scalar in the convolution. The qualitative kernel results of
\cite[Lemma 2.2]{EJC} also apply. In particular, $K$ is even, positive and
integrable, with mass one, and
\begin{equation}\label{eq:kernel-derivative}
 K'(x)=-\frac1{2\sinh(\pi x/2)},\qquad x>0.
\end{equation}
The logarithmic singularity is integrable. All derivatives decay
exponentially away from zero, and
$\int_\R |x|K(x)\dd x<\infty$.

For a physical period $P>0$, set $k=2\pi/P$ and
\begin{equation}\label{eq:periodic-kernel}
 K_P(x)=\sum_{j\in\mathbb Z}K(x+jP).
\end{equation}
The sum converges in $L^1$ on a period and has mass one. Its differentiated
sums converge uniformly on compact sets avoiding $P\mathbb Z$. It is
strictly decreasing and strictly convex on $(0,P/2)$. One way to see the
strict signs for every $P$ is to use
\[
 K(x)=\frac2\pi\sum_{j=0}^\infty
       \frac{e^{-(2j+1)\pi|x|/2}}{2j+1}.
\]
Each periodized exponential is a positive multiple of
$e^{-\mu x}+e^{-\mu(P-x)}$ on $0<x<P$, and has the stated signs on the
half-period.

On writing $\theta=kx$, convolution becomes the even $2\pi$-periodic
Fourier multiplier
\begin{equation}\label{eq:multiplier}
 \K_k\cos(n\theta)=m(kn)\cos(n\theta),\qquad
 m(\xi)=\frac{\tanh\xi}{\xi},\quad m(0)=1.
\end{equation}
For each fixed $k>0$, this symbol has order $-1$. The periodic multiplier
theorem \cite[Proposition 2.7]{EJC} therefore gives
$\K_k:C^s\to C^{s+1}$ at positive noninteger exponents, understood via
the equivalent H\"older--Zygmund spaces. We use the Banach space
\[
 X=C^{1/4}_{\mathrm{even}}(\T),\qquad
 \norm{u}_{C^\alpha}=\norm{u}_\infty+
 \sup_{x\ne y}\frac{|u(x)-u(y)|}{d_\T(x,y)^\alpha}.
\]
Consequently $\K_k:X\to X$ is compact. These mapping bounds may depend
on $P$.

We first record a strict comparison that also applies to continuous
highest profiles.
\begin{lemma}\label{lem:strict-comparison}
If a continuous, even, nonconstant $P$-periodic function $u$ is
nonincreasing on $[0,P/2]$, then $K_P*u$ is strictly decreasing on
$(0,P/2)$. In particular, a solution of $K_P*u=F_c(u)$ taking values in
$[0,\gamma(c)]$ is strictly decreasing there.
\end{lemma}
\begin{proof}
Subtract $\min u$ and write the resulting nonnegative function as the
integral of its superlevel indicators. For every level strictly between
its minimum and maximum, the superlevel set is a centered interval of
radius $r\in(0,P/2)$, up to endpoints. For such an interval,
\[
 \frac{\dd}{\dd x}\bigl(K_P*\mathbf1_{[-r,r]}\bigr)(x)
       =K_P(x+r)-K_P(x-r)<0
\]
for almost every $x\in(0,P/2)$. Indeed the circle distance of $x+r$ is
strictly larger than $|x-r|$, also when $x+r>P/2$.
The interval convolution is locally absolutely continuous, since $K_P$ is
integrable. It is therefore strictly decreasing on every subinterval.
Tonelli's theorem and integration over the nonempty interval of levels
preserve strict decrease. The last assertion follows from the strict
increase of $F_c$ on its stated range.
\end{proof}

\begin{lemma}[A common modulus of continuity]\label{lem:modulus}
Let $1\le c\le\cg$ and let $v$ be a bounded solution of
$K*v=F_c(v)$ on $\R$, with $0\le v\le\gamma(c)$. Periodic solutions
are allowed. With $M=a\cg$, one has
\begin{equation}\label{eq:modulus}
 |v(x+h)-v(x)|^2\le\frac{2M}{\sqrt3}
            \norm{K(\cdot+h)-K}_{L^1(\R)}.
\end{equation}
These profiles have a common $C^\beta(\R)$ bound for every fixed
$0<\beta<1/2$.
\end{lemma}
\begin{proof}
For $0\le s<t\le\gamma(c)$, \eqref{eq:cubic-derivatives} gives
\[
 F_c(t)-F_c(s)\ge\sqrt3c\int_s^t(\gamma(c)-z)\dd z
              \ge\frac{\sqrt3}{2}(t-s)^2.
\]
The last integral, at fixed $t-s$, is smallest when $t=\gamma(c)$.
The equation and $\norm{v}_\infty\le M$ now imply
\eqref{eq:modulus}. Evenness and decrease of $K$ imply, for $h>0$,
\begin{equation}\label{eq:translation-norm}
 \norm{K(\cdot+h)-K}_1=4\int_0^{h/2}K(t)\dd t
                  \le C h(1+|\log h|),\qquad 0<h\le\tfrac12.
\end{equation}
The positive and negative parts change sign at $-h/2$, which gives the
equality. The estimate follows from \eqref{eq:kernel}; the case $h<0$
follows by translation. Combining these bounds proves the claim.
\end{proof}

The estimate requires no crest equality. For each fixed period it gives a
bounded subset of $C^{3/8}$, hence precompactness in $X=C^{1/4}_{\rm even}$.
The latter embedding follows from uniform convergence and the interpolation
estimate
$[f]_{C^\alpha}\le2\norm{f}_\infty d^{-\alpha}
 +[f]_{C^\beta}d^{\beta-\alpha}$, for $\beta>\alpha$.

\begin{lemma}[Speed and small constant states]\label{lem:speed}
A nonconstant positive periodic solution with $v<\gamma(c)$ and $c>0$
has $1<c<\cg$. At $c=1$, the only nonnegative continuous periodic solution
with $0\le v\le\gamma(1)$ is zero. At $c=\cg$, the nonnegative continuous
periodic solutions with $0\le v\le\gamma(\cg)$ are the constants zero and
$\gamma(\cg)$.

For each fixed $P$, all solutions sufficiently near $(v,c)=(0,1)$ in
$X\times\R$ are constant. At every $c_*>1$, zero has a neighborhood
containing no nonconstant solutions.
\end{lemma}
\begin{proof}
Integrating over a period gives $\int(F_c(v)-v)=0$. On
$0<t\le\gamma(c)$ the factor
\[
 \frac{F_c(t)-t}{t}=c^2-1-\frac{3c}{2}t+\frac12t^2
\]
is strictly negative for $0<c\le1$. For $c\ge\cg$ it is nonnegative,
because it decreases up to $\gamma(c)$ and
\begin{equation}\label{eq:upper-speed}
 \frac{F_c(\gamma(c))}{\gamma(c)}=\frac{c^2(1+\sqrt3)}6\ge1.
\end{equation}
At $c=\cg$ equality occurs only at $t=\gamma(c)$ or, before division,
at $t=0$. Continuity gives the constant alternatives.

For the assertion near $(0,1)$, write $v=s+h$, where $s$ is its mean and
$h$ has mean zero. Project $\K_kv-F_c(v)=0$ onto mean-zero functions.
The derivative in $h$ at $(s,h,c)=(0,0,1)$ is $\K_k-I$. Its kernel on
that subspace is zero, since $m(kn)<1$ for all $n\ne0$. As a compact
perturbation of $-I$, it is Fredholm of index zero and is invertible.
The analytic implicit-function theorem, for example
\cite[Lemma 3R]{FNSS}, gives $h=0$ as the unique small solution of the
projected equation for each nearby $(s,c)$. The full equation is thus the
scalar constant equation. Its nearby roots are zero and
\begin{equation}\label{eq:constant-branch}
 b(c)=\frac{3c-\sqrt{c^2+8}}2.
\end{equation}
At zero and $c_*>1$, the full derivative is $\K_k-c_*^2I$.
Convolution has norm at most one for the full H\"older norm, so the
Neumann series gives its inverse. The same implicit-function argument
proves the last assertion. All neighborhoods here may depend on $P$.
\end{proof}

\begin{lemma}[Strict nodal signs]\label{lem:nodal}
A smooth, even, nonconstant periodic solution with $0<v<\gamma(c)$
that is nonincreasing on $[0,P/2]$ satisfies
\[
 v'(x)<0\quad(0<x<P/2),\qquad v''(0)<0,
 \qquad v''(P/2)>0.
\]
This property is open relative to the subcritical solution set in
$X\times\R$.
\end{lemma}
\begin{proof}
The reflected derivative identity is
\begin{equation}\label{eq:reflected-derivative}
 F_c'(v(x))v'(x)=\int_0^{P/2}
       [K_P(x-y)-K_P(x+y)]v'(y)\dd y.
\end{equation}
The bracket is positive inside the half-period. Since $v'\le0$ and is
negative on some interval, the right side is strictly negative. At the
endpoints, differentiation and one integration by parts give
\[
 F_c'(v(0))v''(0)=-2\int_0^{P/2}K_P'(y)v'(y)\dd y<0,
\]
and
\[
 F_c'(v(P/2))v''(P/2)
       =2\int_0^{P/2}K_P'(P/2-y)v'(y)\dd y>0.
\]
The integrals converge because $v'$ vanishes linearly at the endpoints,
whereas $K_P'$ has at most an inverse-distance singularity. This is the
nodal argument of \cite[Theorem 4.2]{EJC}, with physical period $P$.

Near any subcritical solution there is a positive gap from $\gamma(c)$.
The identity $v=F_c^{-1}(\K_kv)$ and the one-derivative gain improve
convergence of solutions in $X$ to convergence in $C^2$. Interior derivative
signs and the endpoint second-derivative signs are then preserved.
\end{proof}

The next estimate prevents the highest waves from becoming flat when the
period grows.
\begin{proposition}[A uniform drop from the critical height]\label{prop:drop}
There are $r_0>0$ and $D_0>0$, independent of $P\ge4$, such that every
smooth nonconstant solution as in Lemma~\ref{lem:nodal}, with
$1<c<\cg$, satisfies
\begin{equation}\label{eq:drop}
 \gamma(c)-v(r)\ge D_0r,\qquad 0<r\le r_0.
\end{equation}
The estimate is inherited by its uniform limits.
\end{proposition}
\begin{proof}
The image terms in $K_P'(t)-K'(t)$ are uniformly bounded for $P\ge4$
and $0<t\le1/2$, by exponential decay. Since $-K'(t)\sim1/(\pi t)$,
we can choose $r_0$ with $2r_0<1/2$ so that
\begin{equation}\label{eq:periodic-derivative-bound}
 -K_P'(2r)\ge\frac1{4\pi r},\qquad 0<r\le r_0,\quad P\ge4.
\end{equation}
On the negative half-period the profile is strictly increasing. For
$x,y\in[-r,-r/2]$, convexity of $K_P$ gives
\[
 K_P(x-y)-K_P(x+y)
   =\int_{|x-y|}^{|x+y|}(-K_P'(t))\dd t
   \ge r[-K_P'(2r)].
\]
The improper nonnegative integral has the same interpretation if the lower
endpoint is zero. Reflecting \eqref{eq:reflected-derivative} to this
half-period, restricting its positive integral to $[-r,-r/2]$, and using
the decrease of $F_c'$ yield
\[
 F_c'(v(-r))v'(x)\ge r[-K_P'(2r)]
                   [v(-r/2)-v(-r)],\qquad -r<x<-r/2.
\]
Integrating in $x$ and cancelling the strictly positive profile increment
gives $F_c'(v(-r))\ge r/(8\pi)$. Put $d=\gamma(c)-v(-r)$. From
\eqref{eq:cubic-derivatives},
\[
 F_c'(v(-r))=\sqrt3c\,d+\frac32d^2
       \le B d,\qquad B=\cg\left(\sqrt3+\frac{3a}{2}\right).
\]
Thus \eqref{eq:drop} holds with $D_0=1/(8\pi B)$. Only smooth
subcritical profiles have been differentiated. The final inequality
passes to uniform limits directly.
\end{proof}

\section{Positive highest periodic waves}\label{sec:periodic}

The positive constant branch \eqref{eq:constant-branch} satisfies
$0<b(c)<\gamma(c)$ for $1<c<\cg$. We construct a nonconstant branch
from it, keeping the physical period fixed throughout this section.

\begin{lemma}[Local bifurcation from the positive constant]\label{lem:local}
For every sufficiently large $P$, there is a unique $c_b\in(1,\cg)$
such that a local analytic branch of even solutions bifurcates from
$(b(c_b),c_b)$ in the first cosine mode. Its positive-amplitude arm has
positive profiles with the nodal signs of Lemma~\ref{lem:nodal}. The
speed is nonconstant, and the profile linearization is invertible at
nearby nonzero points of this branch.
\end{lemma}
\begin{proof}
Set $v=b(c)+w$ and $\Lambda(c)=F_c'(b(c))$. With
$s(c)=\sqrt{c^2/4+2}$, direct calculation gives
\begin{equation}\label{eq:Lambda}
 \Lambda(c)=\frac{c^2}{4}+3-\frac32c s(c),\qquad
 \Lambda'(c)=\frac c2-\frac32\left(s(c)+\frac{c^2}{4s(c)}\right)<0.
\end{equation}
It decreases from $1$ to $0$ on $[1,\cg]$. Since $0<m(k)<1$, there
is exactly one $c_b$ with $\Lambda(c_b)=m(k)$. At that point
$\K_k-\Lambda(c_b)I$ has even kernel spanned by $\cos\theta$.
It is Fredholm of index zero, and its range is precisely the functions
with zero first cosine coefficient. The nonzero derivative in
\eqref{eq:Lambda} gives transversality. The analytic simple-bifurcation
theorem in \cite[Remark 3]{CSV} gives a local curve containing all nearby
nonconstant solutions. Parametrize it by its first cosine coefficient
$\tau$. Half-period translation changes $\tau$ to $-\tau$, so local
uniqueness makes the speed even in $\tau$.

Let $A(c)=3(b(c)-c)/2$. The equation for $w$ is
\begin{equation}\label{eq:w-equation}
 (\K_k-\Lambda(c))w=A(c)w^2+\frac12w^3.
\end{equation}
Analyticity gives expansions, for this fixed $P$,
\[
 w=\tau\cos\theta+\tau^2(w_0+w_2\cos2\theta)+O(\tau^3),
 \qquad c=c_b+C_2\tau^2+O(\tau^4).
\]
All coefficients on the next line are evaluated at $c_b$:
\begin{equation}\label{eq:bifurcation-coefficients}
 \begin{split}
 w_0&=\frac{A}{2(1-m(k))},\qquad
 w_2=\frac{A}{2(m(2k)-m(k))},\\
 -\Lambda'C_2&=A^2\left(\frac1{1-m(k)}
              +\frac1{2(m(2k)-m(k))}\right)+\frac38.
 \end{split}
\end{equation}
The constant and second modes at order two give $w_0,w_2$; the first mode
at order three gives $-\Lambda'C_2=2Aw_0+Aw_2+3/8$.
As $k\downarrow0$,
\[
 1-m(k)=\frac{k^2}{3}+O(k^4),\qquad
 m(k)-m(2k)=k^2+O(k^4),\qquad
 A\longrightarrow-\frac32,\quad \Lambda'\longrightarrow-2.
\]
Consequently
\begin{equation}\label{eq:C2-limit}
 k^2C_2\longrightarrow\frac{45}{16}>0.
\end{equation}
For sufficiently large $P$ this coefficient is nonzero. The local
Lyapunov--Schmidt equation has the form
$G(c,\tau)=\tau(\ell(c-c_b)+q\tau^2+\cdots)$, with
$\ell\ne0$ and $q\ne0$. At its nonzero solutions, differentiation in
$\tau$ with $c$ fixed gives $G_\tau=2q\tau^2+o(\tau^2)$.
This is the profile Schur complement, so the full profile linearization
is invertible there.

The local curve and its expansion can also be constructed in $C^{2+1/4}$;
the same simple kernel and transversality apply, and uniqueness in $X$
identifies the curves. The leading term $b(c)+\tau\cos\theta$ therefore
gives positivity and the strict nodal pattern for $\tau>0$ small.
\end{proof}

We use the corrected analytic global-bifurcation theorem of
Constantin, Strauss and Varvaruca \cite[Theorem 4, Remarks 3--4]{CSV}.
For a real analytic map $\F:\mathcal O\subset\R\times X\to X$, its
hypotheses are: the full trivial line lies in $\mathcal O$; one point on
it has simple kernel and cokernel with transverse parameter crossing;
the profile derivative is Fredholm of index zero at every solution; and
$\mathcal O$ is exhausted by bounded closed sets whose solution subsets
are compact. The distinguished analytic continuation of the local branch
then either eventually leaves every set of that exhaustion or is a
periodic curve. At each branch point, its incoming distinguished analytic
arc has a uniquely selected outgoing arc for which the joined curve admits
a local uniformizing real-analytic parameter
\cite[Remark~3 and Figure~2, preprint pp.~14--15]{CSV}.
This parameter is one-to-one on a sufficiently small interval around the
join. At regular points the corresponding continuation is the unique
implicit-function graph. These are local properties of the selected arc.
We use compact-exhaustion escape,
rather than the stronger norm/boundary disjunction corrected in that paper.

\begin{proposition}[Global continuation and monotonicity]\label{prop:global}
For every sufficiently large fixed $P$, the positive local arm in
Lemma~\ref{lem:local} extends to a distinguished global curve of positive
smooth subcritical profiles. All of its positive-side profiles are
nonconstant and strictly decreasing on $(0,P/2)$. The curve leaves every
compact exhaustion set of the positive subcritical solution domain.
\end{proposition}
\begin{proof}
To place the whole positive constant branch on the trivial line, use the
analytic speed parametrization
\[
 c(\lambda)=1+(\cg-1)\frac{e^\lambda}{1+e^\lambda},\qquad\lambda\in\R,
\]
and define
\[
 \F(\lambda,w)=\K_kw-
  \bigl[F_{c(\lambda)}(b(c(\lambda))+w)-F_{c(\lambda)}(b(c(\lambda)))\bigr]
\]
on the open set
\[
 \mathcal O=\{(\lambda,w)\in\R\times X:
     \min v>0,\ \max v<\gamma(c(\lambda)),\quad v=b(c(\lambda))+w\}.
\]
The whole $w=0$ line belongs to this set. The local crossing is simple
and transverse by Lemma~\ref{lem:local}, since $c'(\lambda)>0$.
At every solution the profile derivative is $\K_k-F_c'(v)$.
Multiplication by $F_c'(v)>0$ is invertible on $X$, whereas $\K_k$ is
compact. The derivative is therefore Fredholm of index zero.

For integer $j\ge1$, take the closed subsets of $\mathcal O$ given by
\[
 |\lambda|,\norm{w}_X\le j,\qquad
 \min v\ge1/j,\qquad \gamma(c)-\max v\ge1/j.
\]
They exhaust $\mathcal O$. On each of them, the inverse cubic has
uniformly bounded derivatives on the relevant range. The equation and
the one-derivative gain give a uniform $C^{1+1/4}$ bound for its
solutions. Their solution subsets are compact in $X$ and closed. This
verifies all hypotheses of the continuation theorem.

It remains to keep the chosen arm in the monotone cone. The cone of even
nonincreasing functions is closed in $X$. At any nonconstant solution in
the cone, Lemma~\ref{lem:nodal} makes the strict nodal property open
relative to the solution curve. Its first loss at a finite parameter
could therefore occur only at a constant. In $\mathcal O$ the only
constant solution is $b(c)$.

At such a contact the profile derivative must be singular, or the
implicit-function theorem would allow only the constant branch nearby.
Its kernel is $\cos(n\theta)$ for exactly one $n\ge1$. The transverse
local bifurcation theorem applies there too. For $n\ge2$, neither sign of
the leading cosine is monotone on $[0,\pi]$, which excludes such a
contact. For $n=1$, strict decrease of $\Lambda$ identifies the contact
with the original bifurcation point.

For completeness, the remaining same-arm return is excluded as follows.
Suppose the first return is at a parameter $T>0$. Arrival along the
negative local arm violates monotonicity already before $T$. Arrival
along the positive local arm identifies its analytic germ with the
outgoing arm near zero, with reversed orientation. Match these two
copies by an order-reversing local parameter map $s\mapsto h(s)$.
The distinguished arc continuation in \cite[Figure 2]{CSV} extends this
match through each regular point and each intervening branch point.
At a regular point this follows from uniqueness of the implicit-function
graph; at a branch point it follows from the unique selection of the
outgoing arc just stated.
Continue while $0<s<h(s)<T$. The left parameter increases and the right
one decreases. If continuation stopped at two distinct interior limits,
the local uniformizing charts at those limits and their transition map
would extend the match.
Thus they meet at some $s_*=h(s_*)$. The two sides at that point would
then describe the identical germ, contrary to local uniformizing
injectivity. This rules out the return. The nearby regular points
needed for the distinguished arc construction are supplied by
\eqref{eq:C2-limit} and the Schur complement in Lemma~\ref{lem:local}.

Hence the entire positive side remains nonconstant and monotone. A
periodic continuation would return to its initial constant point, and is
also excluded. The remaining global alternative is escape from every
compact exhaustion set.
\end{proof}

\begin{proposition}[Highest waves for large periods]\label{prop:highest}
For every sufficiently large $P$, there are $c_P\in(1,\cg)$ and a
positive, even, continuous $P$-periodic solution $v_P$ with
\[
 K*v_P=F_{c_P}(v_P),\qquad v_P(0)=\gamma(c_P).
\]
It is strictly decreasing on $(0,P/2)$, smooth off $P\mathbb Z$, and
satisfies \eqref{eq:drop} and the common modulus
\eqref{eq:modulus}.
\end{proposition}
\begin{proof}
Fix $P$ and follow the positive side of Proposition~\ref{prop:global}.
Lemma~\ref{lem:modulus} makes all its profiles bounded in $C^{3/8}$
and precompact in $X$. Its speeds stay away from $1$: otherwise an
$X$-convergent subsequence has limiting speed $1$, its limiting profile
is zero by Lemma~\ref{lem:speed}, and the constant-only neighborhood
there excludes the nonconstant sequence. Its speeds also stay away from
$\cg$. At that endpoint the limit must be zero or $\gamma(\cg)$.
Local invertibility at zero excludes the first, and the fixed positive
drop \eqref{eq:drop} excludes the second.

The minima stay uniformly positive along this fixed-period curve. Indeed,
a limiting zero minimum at an interior limiting speed would give a
nonnegative solution vanishing at one point. Positivity of $K_P$ then
forces the limit to be identically zero; local invertibility at that
speed again excludes it.

If the crest gap also stayed uniformly positive, the whole positive
curve would lie in one of the exhaustion sets used above. This
contradicts Proposition~\ref{prop:global}. There is therefore a sequence
whose crest gaps tend to zero. Take an $X$-convergent subsequence, also
convergent in speed. Its limit solves the equation and has crest
$\gamma(c_P)$. Estimate \eqref{eq:drop} makes it nonconstant. A zero
would force the limit to vanish identically by positivity of convolution,
so it is positive. It is strictly decreasing by
Lemma~\ref{lem:strict-comparison}. Wherever it is below its crest, the
localized inverse-cubic smoothing of \cite[Lemma 4.1]{EJC}, with kernel
$K_P$, makes it smooth. Strict decrease places these points precisely
off $P\mathbb Z$. Both uniform estimates pass to the limit.
\end{proof}

\section{The solitary limit}\label{sec:solitary}

Choose $P_n\to\infty$ and the highest profiles from
Proposition~\ref{prop:highest}. We keep the physical spatial coordinate
$x$, so the uniform modulus does not change with $n$.

\begin{lemma}\label{lem:limit}
After passage to a subsequence, $c_{P_n}\to c\in[1,\cg]$ and
$v_{P_n}\to v$ locally uniformly on $\R$. The limit is even,
continuous, nonnegative and nonincreasing on $(0,\infty)$, and
\begin{equation}\label{eq:limit-profile}
 K*v=F_c(v),\qquad v(0)=\gamma(c),\qquad
 \gamma(c)-v(r)\ge D_0r\quad(0<r\le r_0).
\end{equation}
In particular, it is nonzero and nonconstant. Its limits at both
infinities agree with a number $b\ge0$ satisfying $F_c(b)=b$.
\end{lemma}
\begin{proof}
The suprema are bounded by $M$ and the modulus in
Lemma~\ref{lem:modulus} is common. Compact-interval selection and a
diagonal subsequence give the stated local uniform convergence; the
speed can be selected first from its compact interval. Evenness and
monotonicity pass to the limit on every compact set. The crest and drop
also pass directly.

For fixed $x$, split the convolution into $|y|\le R$ and $|y|>R$.
On the first part $v_{P_n}(x-y)$ converges uniformly; the complementary
error is bounded by $2M\int_{|y|>R}K(y)\dd y$. Taking $n\to\infty$
and then $R\to\infty$ proves the pointwise equation. The decreasing
bounded positive-half-line profile has a finite limit $b$, also the
negative tail by evenness. Dominated convergence in
$K*v(x)=\int K(y)v(x-y)\dd y$ as $x\to\infty$ yields
$F_c(b)=b$. No global uniform convergence is needed.
\end{proof}

We next show that the tails are zero. The following balance only requires
equal bounded tails, rather than an integrable profile.
\begin{lemma}[First-moment balance]\label{lem:balance}
If $f$ is bounded and has the same finite limit at both infinities, then
\begin{equation}\label{eq:balance}
 \lim_{R\to\infty}\int_{-R}^R(K*f-f)(x)\dd x=0.
\end{equation}
\end{lemma}
\begin{proof}
For each fixed $R$, the absolute double integral for $K*f$ is at most
$2R\norm{f}_\infty\norm{K}_1$. Fubini's theorem and the mass-one
identity give
\[
 \int_{-R}^R(K*f-f)(x)\dd x=\int_\R K(y)A_R(y)\dd y,
 \quad
 A_R(y)=\int_{-R-y}^{R-y}f(x)\dd x-\int_{-R}^Rf(x)\dd x.
\]
For fixed $y$, its two boundary integrals tend to equal and opposite
limits, so $A_R(y)\to0$. For all $R,y$,
$|A_R(y)|\le2\norm{f}_\infty|y|$. The finite first absolute moment
of $K$ supplies an integrable majorant, proving \eqref{eq:balance}.
\end{proof}

\begin{lemma}\label{lem:zero-tail}
The limit in Lemma~\ref{lem:limit} has $1<c<\cg$ and $b=0$.
\end{lemma}
\begin{proof}
If $c=1$, the equation gives
\[
 K*v-v=-\frac12v^2(3-v)\le0.
\]
It is strictly negative on a nonempty open interval, since
$v(0)=a>0$ and $v\le a<3$. Its integrals over expanding symmetric
intervals cannot tend to zero, contrary to Lemma~\ref{lem:balance}.
If $c=\cg$, then $F_c(\gamma(c))=\gamma(c)$. At $x=0$, however,
the positive mass-one average $K*v$ is strictly smaller than
$\gamma(c)$, because $v$ is a nonconstant continuous function bounded
above by $\gamma(c)$. Thus the speed is interior.

Suppose $b>0$. The drop in \eqref{eq:limit-profile} gives
$b<\gamma(c)$. Since $F_c(b)=b$,
\[
 F_c'(b)=1-\frac32cb+b^2<1.
\]
Put $w=v-b\ge0$. It is nonzero, bounded and has zero tails. Strict
concavity of $F_c$ on $[b,\gamma(c)]$ gives
\[
 K*w=F_c(b+w)-F_c(b)<F_c'(b)w<w
\]
wherever $w>0$, with equality in the first difference when $w=0$.
Hence $K*w-w$ is nonpositive everywhere and strictly negative on an
open interval. Lemma~\ref{lem:balance} applied to $w$ is again a
contradiction. This proves $b=0$.
\end{proof}

\begin{lemma}\label{lem:strict-smooth}
The limiting profile is strictly positive on $\R$, strictly decreasing
on $(0,\infty)$, and smooth on $\R\setminus\{0\}$.
\end{lemma}
\begin{proof}
A zero of $v$ would give $K*v>0$ at that point, because $K$ is positive
and $v$ is nonnegative and nonzero. This contradicts $F_c(0)=0$.

For each level $0<t<v(0)$, the superlevel set of $v$ is a centered
interval of finite positive radius $r(t)$, up to endpoints. Its finiteness
uses the zero tail. For every $r>0$, the function
$K*\mathbf1_{[-r,r]}$ is locally absolutely continuous and has derivative
\[
 K(x+r)-K(x-r)<0
\]
for almost every $x>0$. The logarithmic singularity at $x=r$ is
integrable. Each such interval convolution is strictly decreasing.
Integrating over levels by Tonelli's theorem shows that $K*v$ is strictly
decreasing. Strict increase of $F_c$ on $[0,\gamma(c)]$ then gives
strict decrease of $v$. In particular $v(x)<\gamma(c)$ for $x\ne0$.

For smoothness, convolution gains one Sobolev derivative:
$K:H^s(\R)\to H^{s+1}(\R)$, since
$\sqrt{1+\xi^2}\,m(\xi)$ is bounded. On a compact
interval away from zero, localize $v$ with a smooth cutoff equal to one
on a slightly larger interval. The localized part is initially in $L^2$,
so its convolution is in $H^1$. The complementary convolution is smooth
on the smaller interval, since its kernel argument stays away from zero
and every off-origin derivative is integrable at infinity against bounded
$v$. On that interval $F_c'(v)$ has a positive lower bound. Its smooth
scalar inverse, extended off the local range with bounded derivatives,
therefore gives $v\in H^1$ locally by the weak chain rule.

Repeat with nested cutoffs. If $v$ is locally in $H^j$, convolution gains
one derivative. Composition with the smooth local inverse preserves this
integer regularity: in one dimension the lower derivatives in each
chain-rule term are bounded by the local $H^1$ embedding, while at most
one highest derivative is only in $L^2$. Smooth approximation and the
weak product and chain rules justify the calculation. Iteration and
Sobolev embedding yield smoothness off zero.
\end{proof}

\begin{proof}[Proof of Theorem~\ref{thm:main}]
Lemmas~\ref{lem:limit}, \ref{lem:zero-tail} and
\ref{lem:strict-smooth} give all assertions about $c$ and $v$.
Set $\eta=cv-v^2/2$. Expanding the cubic equation yields the exact
identities
\begin{equation}\label{eq:integrated-system}
 K*v+\eta v=c\eta,\qquad \eta+\frac12v^2=cv.
\end{equation}
For any locally integrable profile $f$, the change of variables $z=x-ct$
in a smooth compactly supported space-time test integral gives
$\partial_t f(x-ct)=-c\partial_x f(x-ct)$ in distributions.
Convolution by $K$ commutes with translation since $K\in L^1$ and $v$
is bounded. Applying distributional derivatives to
\eqref{eq:integrated-system} therefore gives both equations of
\eqref{eq:system}. Smoothness makes them classical off the moving crest.
Finally $a-a^2/2=1/3$, so $\eta(0)=c^2/3$.
\end{proof}

\paragraph{AI-use disclosure.}
This work was prepared with extensive assistance from OpenAI Codex,
including mathematical exploration, proof drafting, source inspection and
manuscript preparation. Two additional AI contexts from the same provider
supplied a reader report and a referee report on an anonymous draft; these
reports informed a single revision. Responsibility for the mathematical claims and
references rests with the named author.

\begin{thebibliography}{9}
\bibitem{CSV}
A.~Constantin, W.~Strauss and E.~Varvaruca,
\emph{Global bifurcation of steady gravity water waves with critical layers},
Acta Math. \textbf{217} (2016), 195--262.
\href{https://doi.org/10.1007/s11511-017-0144-x}{doi:10.1007/s11511-017-0144-x}.
The continuation statement cited here is Theorem~4 and Remarks~3--4 in
\href{https://api.newton.ac.uk/website/v0/events/preprints/NI14051}{preprint NI14051}.

\bibitem{EJC}
M.~Ehrnstr\"om, M.~A.~Johnson and K.~M.~Claassen,
\emph{Existence of a highest wave in a fully dispersive two-way shallow
water model}, Arch. Ration. Mech. Anal. \textbf{231} (2019), 1635--1673.
\href{https://doi.org/10.1007/s00205-018-1306-5}{doi:10.1007/s00205-018-1306-5}.
Numbering refers to \href{https://arxiv.org/abs/1610.02603v2}{arXiv:1610.02603v2}.

\bibitem{EMV}
M.~Ehrnstr\"om, O.~I.~H.~M{\ae}hlen and K.~Varholm,
\emph{On the precise cusped behaviour of extreme solutions to Whitham-type
equations}, Ann. Inst. H. Poincar\'e C Anal. Non Lin\'eaire
\textbf{42} (2025), 125--162.
\href{https://doi.org/10.4171/AIHPC/104}{doi:10.4171/AIHPC/104}.

\bibitem{ENW}
M.~Ehrnstr\"om, K.~Nik and C.~Walker,
\emph{A direct construction of a full family of Whitham solitary waves},
\href{https://arxiv.org/abs/2204.03274v1}{arXiv:2204.03274v1} (2022).

\bibitem{FNSS}
S.~Fu\v c\'ik, J.~Ne\v cas, J.~Sou\v cek and V.~Sou\v cek,
\emph{Upper bound for the number of eigenvalues for nonlinear operators},
Ann. Scuola Norm. Sup. Pisa Cl. Sci. (3) \textbf{27} (1973), 53--71.
\href{https://www.numdam.org/item/ASNSP_1973_3_27_1_53_0/}{Primary archive}.

\bibitem{NW}
D.~Nilsson and Y.~Wang,
\emph{Solitary wave solutions to a class of Whitham--Boussinesq systems},
Z. Angew. Math. Phys. \textbf{70} (2019), article 70.
\href{https://arxiv.org/abs/1810.03405v2}{arXiv:1810.03405v2}.
\end{thebibliography}
\end{document}
